\section{The exact problem and its reductions}\label{sec:exact}

Let $a^2+b^2=c^2$ with positive integers, set
$(m_1,m_2,m_3)=(a,b,c)$, and put
\[
q_1=(-c/2,0),\quad q_2=(c/2,0),\quad
q_3=((b^2-a^2)/(2c),ab/c),\qquad \dot q_i=0.
\]
The side opposite each mass has the same length as that mass.  We work with $G=1$ and the equations
\[
 \ddot q_i=\sum_{j\ne i}m_j{q_j-q_i\over|q_j-q_i|^3}.
\]
\begin{conjecture}[Pythagorean--Burrau nonperiodicity]
For every positive integer triple $a^2+b^2=c^2$, the solution with these
initial data has no collision-free labelled period.
\end{conjecture}
A labelled period is an exact return in the fixed inertial frame. The
maximal classical solution ends at its first collision; regularized
continuation through that collision is a different notion of solution.

\subsection{Scaling and the rational parameter curve}

\begin{lemma}[Simultaneous mass--length scaling]
If $q_i(t)$ solves the equations with masses $m_i$, then for $k>0$
\[
\widetilde m_i=km_i,\qquad \widetilde q_i(t)=kq_i(t/k)
\]
solves the rescaled problem.  A labelled period $T$ corresponds exactly to a
period $kT$.
\end{lemma}
\begin{proof}
Both sides of Newton's acceleration equation acquire the factor $k^{-1}$.
Velocity is unchanged after the corresponding time substitution, and mutual
distances acquire the positive factor $k$.
\end{proof}

It follows that only primitive triples matter.  For $u\in(0,1)$ define
\[
A(u)=\frac{1-u^2}{1+u^2},\qquad
B(u)=\frac{2u}{1+u^2}.
\]
If $u=p/q$ is reduced, Euclid's integer triple is
$(q^2-p^2,2pq,q^2+p^2)$ and its common divisor is two exactly when $p,q$ are
both odd.  Conversely Euclid's theorem recovers every ordered primitive
triple.  The inverse is $u=B/(1+A)$, and leg exchange is the involution
$u\mapsto(1-u)/(1+u)$.  Its fixed point is $\sqrt2-1$, giving the fundamental
interval $0<u\leq\sqrt2-1$ for the real family.

After normalizing $c=1$, masses and positions are
\[
(A,B,1),\qquad q_1=(-1/2,0),\quad q_2=(1/2,0),\quad
q_3=((B^2-A^2)/2,AB).
\]
We call this one-parameter family \emph{tied}, since its masses and initial
opposite sides vary together. The stronger real conjecture makes the same
nonperiodicity assertion for every $0<u\le\sqrt2-1$.

\subsection{Initial geometry}

The initial center of mass is the incenter: both have barycentric weights
equal to the opposite side lengths. The identity
$I=(\sum m_i)^{-1}\sum_{i<j}m_im_jr_{ij}^2$ gives, directly,
\[
 I_0=abc,\qquad U_0={ab\over c}+{ac\over b}+{bc\over a}.
\]
In the normalized problem,
\[
 I_0=AB,\qquad U_0=AB+(AB)^{-1},\qquad H=-U_0.
\]
Here $I$ is the moment of inertia about the center of mass and $U$ is the
positive potential magnitude. The right angle is not preserved: for
$D=r_{12}^2-r_{23}^2-r_{31}^2$,
\[
 D(0)=\dot D(0)=0,\qquad
 \ddot D(0)=2\bigl[(AB)^{-1}-AB(A+B)\bigr]>0.
\]
Thus the angle opposite $r_{12}$ initially becomes obtuse. This local sign
does not define an invariant region; numerical integration at $u=1/3$
finds later changes of sign of $D$.

\subsection{Second brakes and a residual valid at collinearity}

\begin{lemma}[Second-brake lemma]
A classical free-fall segment generates a collision-free labelled periodic
solution if and only if all inertial velocities vanish at some positive time
before collision.
\end{lemma}
\begin{proof}
Uniqueness and time reversal give $q(-t)=q(t)$.  A brake at $\tau$ gives
$q(\tau+s)=q(\tau-s)$ and hence a labelled period $2\tau$.  Conversely, a
labelled period $T$ combined with the symmetry at zero gives reflection
symmetry about $T/2$, and therefore zero velocity there.  The period $2\tau$
need not be minimal unless $\tau$ is the first positive brake.
\end{proof}

Write $r_{ij}=|q_j-q_i|$, $K=\tfrac12\sum_i m_i|\dot q_i|^2$,
$U=\sum_{i<j}m_im_j/r_{ij}$, and $H=K-U=-U_0$. In the center-of-mass
frame define reduced masses $\mu_1=m_1m_2/(m_1+m_2)$ and
$\mu_2=m_3(m_1+m_2)/(m_1+m_2+m_3)$, and Jacobi vectors
\[
X=q_2-q_1,\qquad
Y=q_3-\frac{m_1q_1+m_2q_2}{m_1+m_2}
\]
and Hopf invariants
\[
h(X,Y)=\left(\frac{|X|^2-|Y|^2}{2},X\cdot Y,X\times Y\right).
\]
We use the quotient residual
\[
\Bcal(u,t)=\frac{d}{dt}h(X(u,t),Y(u,t)).
\]
Away from triple collision, $Dh$ has rank three and kernel equal to the common
rotation direction.  Since total linear and angular momentum vanish,
$\Bcal=0$ is equivalent to vanishing of every labelled inertial velocity.
Unlike three mutual-distance derivatives, this remains true at syzygy.
Indeed, a velocity in $\ker Dh$ has the form
$(\dot X,\dot Y)=\omega(iX,iY)$, whose angular momentum is $\omega I$;
$L=0$ and $I>0$ force $\omega=0$.

For an explicit configuration-space reduction, rotate with $X$ and write
$X=(R,0)$, $Y=(C,D)$. Eliminating the common rotation at $L=0$ gives
\[
 K_{\rm red}={\mu_1\dot R^2+\mu_2(\dot C^2+\dot D^2)\over2}
 -{\mu_2^2(C\dot D-D\dot C)^2\over
 2[\mu_1R^2+\mu_2(C^2+D^2)]}.
\]
Here the dots on $C,D$ differentiate coordinates in the moving frame.
Substitution of
$\dot\theta=-\mu_2(C\dot D-D\dot C)/I$ into the inertial kinetic
energy proves the formula. With $m_{12}=m_1+m_2$, the potential is
\[
 U={m_1m_2\over R}
 +{m_1m_3\over\sqrt{(C+m_2R/m_{12})^2+D^2}}
 +{m_2m_3\over\sqrt{(C-m_1R/m_{12})^2+D^2}}.
\]
The tied brake curve is
$R=1$, $C=AB(B-A)/(A+B)$, $D=AB$, with zero reduced velocity.

\section{Finite certificates for nonperiodicity}\label{sec:events}

Let
\[
 I=\mu_1|X|^2+\mu_2|Y|^2,\qquad
 \sigma=\mu_1\overline X\dot X+\mu_2\overline Y\dot Y,
 \qquad
 \zeta=\mu_1\overline X\dot X-\mu_2\overline Y\dot Y .
\]
Here complex multiplication identifies the oriented plane with $\mathbb C$.
Then $\dot I=2\operatorname{Re}\sigma$ and the angular momentum is
$L=\operatorname{Im}\sigma$.  Consequently, at $L=0$ and away from $Y=0$,
\[
 \dot I=0\ \hbox{ and }\ \zeta=0
 \quad\Longleftrightarrow\quad
 \dot X=\dot Y=0.
\]
The implication from a brake to these three scalar equalities is
unconditional.  At the Jacobi degeneracy $Y=0$, the Hopf residual $\Bcal$
above supplies the coordinate-independent converse.

This splitting has an exact shape interpretation.  Put
$J_X=\mu_1|X|^2$, $J_Y=\mu_2|Y|^2$, $s=J_X/I$, and
$\phi=\arg Y-\arg X$.  At every zero of $\dot I$ with $X,Y\ne0$ and $L=0$,
\[
 \boxed{\quad
 \zeta=I\left(\dot s-2is(1-s)\dot\phi\right),\qquad
 |\zeta|^2={8KJ_XJ_Y\over I}=8KI\,s(1-s).
 \quad}
\]
Indeed $\sigma=0$ gives
$\zeta=2\mu_1\overline X\dot X=-2\mu_2\overline Y\dot Y$; decomposing each
Jacobi velocity into radial and angular parts proves both formulas.  Thus a
maximum event is a brake exactly when its shape velocity vanishes, and the
origin-avoidance problem can equivalently be written as the scalar strict
inequality $K=U-U_0>0$. At a brake time $\tau$, time reversal gives
\[
 K(\tau+s)={s^2\over2}\sum_i m_i|\ddot q_i(\tau)|^2+O(s^4).
\]
The coefficient is positive, since $\ddot I(\tau)=-2U_0<0$ rules out
vanishing of all accelerations. A brake is consequently a quadratic zero
of $K$ in time, which cannot be detected by a sign change of $K$.

There is also a brake residual intrinsic to any binary Levi--Civita chart.
For a selected pair write $g=q_j-q_i=w^2$, $dt=|w|^2d\sigma$, and
$z=dw/d\sigma$. Let $G$ be the third body relative to the selected pair's
center of mass and $P=\dot G$. With reduced masses $\mu_g,\mu_G$, angular
momentum is
\[
 L=2\mu_g\operatorname{Im}(\overline w z)+\mu_G G\times P.
\]
Thus, on a collision-free zero-angular-momentum trajectory and wherever
$G\ne0$,
\[
 \boxed{\quad (z_r,z_i,G\cdot P)=0
 \quad\Longleftrightarrow\quad\hbox{a labelled brake}.\quad}
\]
Indeed $z=0$ gives $\dot g=0$; then $L=0$ gives $G\times P=0$, while
$G\cdot P=0$ forces $P=0$. Total momentum removes the common translation.
This residual stays regular through arbitrarily close selected-pair
passages; at $G=0$ one changes Jacobi tree or uses $\Bcal$.

\begin{theorem}[Brake-event identities]
Along a collision-free free-fall solution with initial potential $U_0$, every
brake satisfies
\[
 K=0,\qquad U=U_0,\qquad \ddot I=-2U_0<0,
 \qquad r_{ij}\geq {m_im_j\over U_0}.
\]
At any zero of $\dot I$ for which $\ddot I\geq0$, one instead has
$K\geq U_0$.
\end{theorem}
\begin{proof}
With positive potential magnitude $U$ and energy $H=K-U=-U_0$, the
Lagrange--Jacobi identity is
\[
 \ddot I=4K-2U=2K-2U_0=2U-4U_0.
\]
At a brake $K=0$, giving the first three assertions.  Since every summand
$m_im_j/r_{ij}$ is positive and bounded above by $U=U_0$, the separation
bound follows.  If $\ddot I\geq0$, the middle expression instead gives
$K\geq U_0$.
\end{proof}

\begin{corollary}[Validated-cover criterion]\label{cor:cover}
Suppose a classical free-fall solution has a launch window $[0,t_L]$ on
which $U<2U_0$. Suppose validated flow enclosures cover every time in
$[t_L,t_*]$, and every enclosure proves at least one of
\[
 \dot I\ne0,\qquad K>0,\qquad \Bcal\ne0.
\]
Then there is no second brake in $(0,t_*]$. If the terminal enclosure
satisfies a proved forward collision-or-escape criterion that excludes every
later brake, the initial free-fall solution is nonperiodic.
\end{corollary}
\begin{proof}
At launch $\dot I(0)=0$, whereas $U<2U_0$ implies $\ddot I<0$.
Consequently $\dot I(t)<0$ for $0<t\le t_L$. The remaining disjunction
excludes a brake at every covered time. This argument applies fiber by fiber
when $t_*$ depends on a parameter; both section images and the complete
intervening flow tubes must be covered. The initial brake itself is never
subjected to the nonvanishing disjunction.
\end{proof}

This criterion has been implemented with exact rational initial data and
MPFR interval arithmetic.  The archived certificates prove the conjecture
for
\[
(21,20,29),\ (5311,5280,7489),\ (8319,8200,11681),
\ (33439,32400,46561),\ (72,65,97)
\]
in both leg orderings. The parameter-cover results below require additional
control of how the enclosed trajectories depend on $u$.

\begin{lemma}[Terminal escape certificate]\label{lem:escape}
For a prospective binary of total mass $M$ and an outer body of mass $m_c$,
let $x,y$ be the inner and outer Jacobi vectors, $r=|x|$, $\rho=|y|$, and
$e=|\dot x|^2/2-M/r$. Choose $\eta>0$, set $R=M/\eta$ and
$s_0=\rho_0-R$, and suppose
\[
s_0>0,\quad \dot\rho_0>0,\quad
\delta=\frac12\dot\rho_0^2-\frac{M+m_c}{s_0}>0,
\]
\[
e_0+\frac{m_c\sqrt{2MR}}
{\sqrt{2\delta}\,s_0^2}<-\eta.
\]
Then the future solution either ends in inner binary collision or the outer
body escapes with $\dot\rho\ge\sqrt{2\delta}$ forever. In particular there is
no later classical brake.
\end{lemma}
\begin{proof}
On the bootstrap region $e<-\eta$, one has
$r<R$ and $r|\dot x|<\sqrt{2MR}$. The Newton field derivative bound
$\|D(v/|v|^3)\|\le2/|v|^3$ gives
\[
|\dot e|\le\frac{2m_cr|\dot x|}{(\rho-r)^3}.
\]
The outer Jacobi acceleration obeys
$|\ddot y|\le(M+m_c)/(\rho-r)^2$, so
$\ddot\rho\ge-(M+m_c)/(\rho-R)^2$. Therefore
\[
J=\frac12\dot\rho^2-\frac{M+m_c}{\rho-R}
\]
is nondecreasing while $\dot\rho>0$. Its initial value $\delta>0$
prevents $\dot\rho$ from reaching zero. Hence
$\rho\ge\rho_0+\sqrt{2\delta}(t-t_0)$, and the total future tidal increase of
$e$ is at most the allowance in the hypothesis. Its strict margin closes the
bootstrap. The positive gap $\rho-r$ excludes an outer collision; only an
inner binary collision can terminate the resulting escape trajectory.
\end{proof}

The near-triple-collision scaling below has limiting inner energy zero,
so it requires an escape criterion that permits an expanding inner pair.

\begin{lemma}[Escape with a possibly expanding inner pair]\label{lem:hierarchical}
Use the same Jacobi variables, let $w_0=\rho_0-r_0>0$, and choose
$\varepsilon,c>0$. Put
\[
 v_b=\sqrt{2M/r_0+2\varepsilon},\qquad
 \mathcal T=\int_0^\infty
 {\sqrt{2M(r_0+v_bt)+2\varepsilon(r_0+v_bt)^2}
  \over(w_0+ct)^3}\,dt.
\]
If
\[
 \dot\rho_0-{M+m_c\over cw_0}>v_b+c,\qquad
 e_0+2m_c\mathcal T<\varepsilon,
\]
then, until a possible inner collision,
$r(t)\le r_0+v_bt$, $\dot\rho(t)>v_b+c$, and
$\rho(t)-r(t)\ge w_0+ct$. In particular there is no later brake.
\end{lemma}
\begin{proof}
Assume $e<\varepsilon$ and $\rho-r\ge w_0+ct$ on a bootstrap interval.
At a first outward crossing of $r_0+v_bt$, the energy inequality gives
$\dot r\le|\dot x|<v_b$, excluding that crossing. Integration of
$\ddot\rho\ge-(M+m_c)/(w_0+ct)^2$ then gives
$\dot\rho>v_b+c$ and preserves the gap. Finally
$r|\dot x|\le\sqrt{2Mr+2\varepsilon r^2}$ bounds the total tidal
increase of $e$ by $2m_c\mathcal T$. The integral converges, and the
strict energy margin closes the bootstrap.
\end{proof}

\begin{theorem}[Validated middle first-maximum interval]
For every real
\[
 {29\over100}\le u\le {14501\over50000},
\]
the tied classical trajectory is collision-free through its first positive
local maximum of $I$, and this maximum is not a labelled brake.  The unique
intervening minimum and the first subsequent maximum obey
\[
 t_{\min}\in[0.758038,0.758045],\qquad
 t_{\max}\in[1.32908,1.33943].
\]
\end{theorem}
\begin{proof}
The exact tied launch is embedded as a one-dimensional mean-value graph with
$u$ retained as a frozen state coordinate.  A pinned MPFR/CAPD $C^1$
Poincare calculation propagates this graph first in a pair--13
Levi--Civita chart and, after an exact Jacobi-tree transformation at $t=1$,
in a pair--23 chart.  On every accepted-step enclosure, all three mutual
separations are strictly positive.  The launch-concavity window and whole-leg audits give $\dot I<0$ after
launch to the first minus-to-plus zero and $\dot I>0$ from that zero to the
first later plus-to-minus zero.  At the minimum,
$U/U_0\in[5.17942,5.18335]$; at the maximum,
$U/U_0\in[0.983016,1.01907]\subset(0,2)$, proving both events strict by
Lagrange--Jacobi.  In the final selected-pair chart the validated residual
has
\[
 \operatorname{Im}z\in[-0.0500738,-0.0113187],\qquad
 P_x\in[-0.0561842,-0.00917817].
\]
A labelled brake forces $z=P=0$, so either strict component excludes it.
Two abutting tiles cover the stated interval. The exact chart identities,
logs, and an independent left-tile replay are indexed in
Appendix~\ref{sec:reproduction}.
\end{proof}

This theorem excludes only the first positive maximum.  It does not rule out
a brake on a later maximum branch and therefore is not, by itself, a
nonperiodicity theorem on the displayed interval.

\begin{theorem}[Validated middle real interval]
For every real
\[
 {29\over100}\le u\le {2900000101\over10000000000},
\]
the tied maximal classical solution has no second labelled brake at any
collision-free time and hence is not a classical periodic solution.
Consequently the Pythagorean--Burrau conjecture holds for infinitely many
primitive triples in this interval, including
$(9159,5800,10841)$ at $u=29/100$.
\end{theorem}

More explicitly, for every integer $k\ge99010$ set
\[
 u_k={290k+1\over1000k}={29\over100}+{1\over1000k}.
\]
The numerator and denominator are coprime and of opposite parity, and the
interval inequality is equivalent to $1/(1000k)\le101/10^{10}$.  The theorem
therefore proves nonperiodicity for the primitive triples
\[
 \left(915900k^2-580k-1,
       580000k^2+2000k,
       1084100k^2+580k+1\right),
 \qquad k\ge99010.
\]
\begin{proof}
A pinned MPFR/CAPD propagation embeds the exact launch curve on
$[0.29,0.290000001]$ in one correlated tripleton with $u$ frozen as a state
coordinate.  Before the
exchange, every accepted full-step enclosure is checked for three positive
mutual separations. The launch window verifies $U<2U_0$; subsequent windows
verify the brake-exclusion disjunction of Corollary~\ref{cor:cover}.  Three algebraic changes of Levi--Civita chart use rigorous
mean-value images.

Across the exchange, six $C^1$ Poincare maps project onto the transverse
pair--13 sections
\[
 w_r=-{3\over5},-{1\over2},-{2\over5},-{3\over10},-{1\over5},0.
\]
For each map a convex mean-value domain contains the propagated family and
its stored center, the validated section derivative transports the
distinguished parameter generator, and an independent $C^0$ integration
covers the complete incoming tube through the latest possible return.
Thus section reconditioning does not omit any inter-section time.  Every
tube enclosure retains positive separation and excludes a brake.

For the terminal binary $\{2,3\}$ set $M=B+1$, $G=q_1-C_{23}$,
$\rho=|G|$, $h=|\dot q_3-\dot q_2|^2/2-M/r_{23}$, and
\[
 R=M/4,\quad d=\rho-R,\quad
 E_\rho={\dot\rho^2\over2}-{A+B+1\over d},\quad
 \Delta={A\sqrt{2MR}\over\sqrt{2E_\rho}\,d^2}.
\]
These are precisely the quantities of Lemma~\ref{lem:escape}; $h$ is
transported as a regularized state variable. Hence the terminal check can
use its enclosure without dividing by the small inner separation.
The resulting terminal enclosure has
\[
\begin{split}
t_*&\in[4.30030837364,4.30032429905],\\
d&>1.8655,\qquad \dot\rho>2.0484,\qquad E_\rho>0.8223,\\
h&<-9.2745,\qquad -4-h-\Delta>5.0690 .
\end{split}
\]
These are the strict hypotheses of Lemma~\ref{lem:escape}
with binary $\{2,3\}$, escaper body 1, and $\eta=4$.  Hence after $t_*$ the
inner pair either collides classically, ending the maximal classical
solution, or body 1 escapes permanently; neither alternative permits a later
brake. The second-brake lemma proves nonperiodicity. A separate interval
checker verifies the terminal inequalities on an outward-widened box;
an independent MPFR replay reproduces the section audits and escape margin.

The remaining extension is the union of 100 independently certified,
adjacent closed tiles of exact width $10^{-10}$, beginning at
$0.2900000001$ and ending at $0.2900000101$.  Every requested interval is
identical to its PASS interval, every consecutive pair shares its exact
rational endpoint, and all 100 full logs are archived. The campaign overlaps
the initial tile, and their union is exactly the stated closed interval.
\end{proof}

This continuum theorem covers a width of $1.01\times10^{-8}$. Larger
parameter covers remain a separate validation problem; experimental wider
propagations are not included in its conclusion.

At a possible brake, the bound $r_{ij}\ge m_im_j/U_0$ is uniform on every
compact parameter interval bounded away from $u=0$. It gives no lower
bound on separations between brake events and degenerates as $u\to0$.

On any regular strict-maximum branch, the implicit-function theorem locally
gives $t=t_k(u)$.  The unresolved condition becomes the planar
origin-avoidance problem
\[
 u\longmapsto \zeta(u,t_k(u))\in\mathbb C\setminus\{0\}.
\]
To turn this into a universal theorem one must also control branch creation,
collision and escape endpoints, degenerate events, and the possible number
of branches.  The nominal dimension count alone supplies none of these
facts.

Numerical diagnostics at $u=1/3$ place maximum-event values of $\zeta$ in
all four quadrants, with the origin in the convex hull of three early
values. These data argue against a fixed real linear projection with one
strict sign at every maximum. Branch-dependent order or a nonlinear
constraint remains a possible approach.
